\begin{textsample}{NEG Dim 1 – 2009\_c6 – Score -3.00 – 2009\_c6/enade\_2009\_tec\_mark\_8016f58348.txt}  \label{ex:f1_neg_214}
Capítulo 3 Análise Técnica da Prova Este capítulo tem por objetivo apresentar o desempenho dos estudantes de TECNOLOGIA EM MARKETING no ENADE/2009.
Para isso, foram calculadas as estatísticas básicas da prova como um todo, bem como as estatísticas das questões relacionadas à Formação Geral e ao Componente Específico.
Nas Tabelas, são apresentadas as seguintes estatísticas : total da população, de presentes, média, erro-padrão da média, desvio-padrão, mediana, nota máxima, nota mínima e coeficiente de assimetria ( ver definição em 1.4.3 ).
Nas estatísticas das questões discursivas, foram detalhadas também as notas nulas : em branco, por protesto ou por insuficiência.
As estatísticas apresentadas neste capítulo contemplaram, separadamente, os ingressantes, os concluintes e o total de estudantes, e foram calculadas tendo -se em vista as seguintes agregações : ( a ) Região Geográfica e País e ( b ) a Categoria Administrativa e a Organização Acadêmica.
As notas brutas das questões discursivas em \textbf{cada} uma das partes, Formação Geral/Componente Específico, foram calculadas a partir da média simples das notas de \textbf{cada} uma das questões correspondentes.
A primeira questão discursiva de Formação Geral foi desconsiderada e, portanto, a nota bruta referente às questões discursivas de Formação Geral foi igual à nota da questão mantida, a segunda.
No Componente Específico, a nota bruta correspondeu à média das notas nas três questões.
As notas brutas das questões objetivas de \textbf{cada} uma das partes, Formação Geral/Componente Específico, foram calculadas a partir dos acertos dos estudantes considerando -se uma regra de três para a conversão para a escala de 0 a 100, descontando -se aquelas questões que foram desconsideradas pela comissão ou que não passaram pelo teste do índice de discriminação ponto-bisserial ( ver seção 1.4.6 para a descrição do índice ).
As questões com índices de discriminação inferior a 0,2 foram excluídas do cômputo da nota do estudante.
As notas brutas de Formação Geral e do Componente Específico para \textbf{cada} um dos estudantes foram calculadas como médias ponderadas das suas notas nas questões objetivas e discursivas, ponderação diferente para \textbf{cada} parte : C\_FG = 0,6 × C\_FG^OBJ + 0,4 × C\_FG^DIS C\_CE = 0,85 × C\_CE^OBJ + 0,15 × C\_CE^DIS, 3.1 Estatísticas Básicas da Prova 3.1.1 Estatísticas Básicas Gerais A Tabela 3.1 apresenta as estatísticas básicas da prova por grupo de estudantes.
O total de estudantes inscritos para a prova foi de 23.301.
Destes, 67,0 \% compareceram, sendo que a presença foi maior entre concluintes ( 78,1 \% ) do que entre ingressantes ( 59,5 \% ).
A média geral da prova foi de 48,00, sendo que os ingressantes obtiveram média mais baixa, de 45,20, que os concluintes, cuja média foi 51,20.
O desvio-padrão geral foi de 15,6, sendo o do grupo de ingressantes menor ( 15,2 ), do que o do grupo dos concluintes ( 15,4 ), aqueles possuíram uma distribuição mais homogênea nas notas.
A nota máxima foi 95,7, obtida por pelo menos um concluinte ao passo que a maior nota obtida por um ingressante foi 91,4.
Os coeficientes de assimetria do total de estudantes e do grupo de ingressantes foram simétricos, no entanto o do grupo de concluintes foi negativo, indicando neste uma maior concentração à direita da distribuição e um maior espalhamento à esquerda. onde C\_FG é a nota bruta de Formação Geral, C\_FG^OBJ é a nota bruta nas questões objetivas de Formação Geral, C\_FG^DIS é a nota bruta nas questões discursivas de Formação Geral, C\_CE é a nota bruta do Componente Específico, C\_CE^OBJ é a nota bruta nas questões objetivas do Componente Específico, e C\_CE^DIS é a nota bruta nas questões discursivas do Componente Específico.
A nota bruta da prova, C\_GER, por sua vez, é também uma média ponderada da nota de \textbf{cada} uma das partes, Formação Geral/Componente Específico : C\_GER = C\_FG × 0,25 + C\_CE × 0,75. 3.1 Estatísticas Básicas da Prova 3.1.1 Estatísticas Básicas Gerais A Tabela 3.1 apresenta as estatísticas básicas da prova por grupo de estudantes.
O total de estudantes inscritos para a prova foi de 23.301.
Destes, 67,0 \% compareceram, sendo que a presença foi maior entre concluintes ( 78,1 \% ) do que entre ingressantes ( 59,5 \% ).
A média geral da prova foi de 48,00, sendo que os ingressantes obtiveram média mais baixa, de 45,20, que os concluintes, cuja média foi 51,20.
O desvio-padrão geral foi de 15,6, sendo o do grupo de ingressantes menor ( 15,2 ), do que o do grupo dos concluintes ( 15,4 ), aqueles possuíram uma distribuição mais homogênea nas notas.
A nota máxima foi 95,7, obtida por pelo menos um concluinte ao passo que a maior nota obtida por um ingressante foi 91,4.
Os coeficientes de assimetria do total de estudantes e do grupo de ingressantes foram simétricos, no entanto o do grupo de concluintes foi negativo, indicando neste uma maior concentração à direita da distribuição e um maior espalhamento à esquerda.
Uma comparação entre o desempenho de ingressantes e concluintes pode ser realizada a partir dos dados apresentados no Gráfico 3.1.
Cumpre destacar que as notas dos concluintes, de um modo geral, foram mais elevadas do que as dos ingressantes, cuja distribuição de notas \textbf{está} à esquerda dos primeiros, superioridade reforçada pelas notas medianas dos dois grupos : 51,6 entre os concluintes e 45,1 entre os ingressantes.
Nas faixas de notas até 50 pontos o percentual de ingressantes é maior que o dos concluintes, situação esta que se inverte para as faixas de notas superiores a 50 pontos.
Constata -se que aproximadamente 38,0 \% dos ingressantes e 54,0 \% dos concluintes situaram -se nas faixas de notas superiores a 50,0 pontos.
O Gráfico 3.2 apresenta informações referentes ao desempenho geral de ingressantes e concluintes, comparando os resultados considerando a Região, categoria administrativa e Organização Acadêmica.
Levando -se em conta as notas médias dos estudantes em \textbf{cada} Região, observa -se que, nos dois grupos há uma variação expressiva entre as Regiões.
As notas médias mais elevadas entre os concluintes ocorreram nas Regiões Sul ( 51,98 ) e Sudeste ( 51,79 ), ambas acima da média nacional ( 51,20 ).
Em relação aos ingressantes, as médias mais altas foram obtidas nas Regiões Sul ( 47,85 ) e Sudeste ( 45,56 ).
Nas demais Regiões a média foi inferior à média nacional ( 45,19 ).
Na Região Centro-Oeste foi obtido a média de notas mais baixa entre os concluintes ( 44,68 ) e ingressantes ( 41,16 ).
Cursos nesta área só foram oferecidos em instituições Municipais e Particulares.
Levando -se em conta os agrupamentos dos estudantes em categorias administrativas, observa -se que a pontuação média mais elevada entre os concluintes e ingressantes foi encontrada em instituições de origem municipal ( respectivamente 51,75 e 46,73 ) – ambas acima da média nacional.
Tendo como foco as notas médias dos estudantes agrupados por Organização Acadêmica, notamos que as faculdades apresentam as maiores médias entre os concluintes ( 51,43 ) e os centros universitários apresentam as maiores médias entre os ingressantes ( 46,16 ), ao passo que as piores notas médias foram observadas nas universidades para ambos os grupos. 3.1.2 Estatísticas Básicas de Formação Geral A Tabela 3.2 apresenta as estatísticas básicas de Formação Geral por grupo de estudantes.
O total de estudantes inscritos para a prova foi de 23.301.
Destes, 67,0 \% compareceram, sendo que a presença foi maior entre concluintes ( 78,1 \% ) do que entre ingressantes ( 59,5 \% ).
A média geral da parte da prova de Formação Geral foi de 39,15, sendo que os ingressantes obtiveram média mais baixa, 37,55, que os concluintes, cuja média foi 40,98.
O desvio-padrão geral foi de 20,4, sendo o mesmo para ingressantes e concluintes.
A nota máxima foi 100,0, obtida por estudantes de ambos os grupos.
Para este componente, assim como para a nota geral, os coeficientes de assimetria são positivos para todos os grupos.
O Gráfico 3.3 permite a comparação entre concluintes e ingressantes relativamente ao desempenho na parte da prova que aborda a Formação Geral.
Nesta parte, se excluindo as notas nulas, tal como na prova como um todo, também os concluintes apresentaram desempenho superior, embora seja perceptível um maior equilíbrio.
A mediana dos concluintes foi superior à dos ingressantes, o mesmo acontecendo com a média.
Dos ingressantes, 26,4 \% alcançaram nota superior a 50 ; dentre os concluintes, o percentual dos estudantes nesta faixa de notas foi 32,5 \%.
No Gráfico 3.4 são apresentadas as informações referentes ao desempenho dos ingressantes e dos concluintes na parte de Formação Geral, em diferentes agregações como Região, categoria administrativa e Organização Acadêmica.
Considerando -se as notas médias dos estudantes segundo \textbf{cada} Região, observou -se que, tanto no grupo dos ingressantes quanto no grupo dos concluintes, houve diferenças marcantes.
As médias variam, entre os concluintes, de 33,14 na Região Norte a 41,94 na Região Sul.
Entre os ingressantes, os valores foram de 32,16 na Região Norte a 40,37 na Região Sul.
A maior dispersão de notas dos estudantes de uma dada Região foi encontrada no Norte, tanto para os concluintes quanto para os ingressantes.
Em todas as Regiões os concluintes apresentaram médias maiores do que os ingressantes.
Todas as Regiões apresentaram concluintes.
Tendo como foco as categorias administrativas, destacou -se que somente as instituições municipais e particulares \textbf{estavam} representadas.
Sendo que as notas médias destas últimas foram superiores para concluintes ( 40,98 ) às das instituições municipais ( 40,01 ).
Por outro lado, as notas médias das particulares foram inferiores para ingressantes ( 37,50 ), às das instituições municipais ( 43,32 ).
Considerando -se o tipo de Organização Acadêmica, os Centros Universitários apresentaram os maiores valores, tanto para concluintes ( 42,02 ) quanto para ingressantes ( 38,79 ).
Os Centros Federais de Educação Tecnológica e os Institutos Federais de Educação, Ciência e Tecnologia não \textbf{estavam} representadas nesta área. 3.1.3 Estatísticas Básicas do Componente Específico A Tabela 3.3 apresenta as estatísticas básicas referentes ao Componente Específico da prova, tendo sua média geral, de 51,00, superior ao desempenho da prova de TECNOLOGIA EM MARKETING como um todo, cuja média foi 48,00, como apresentado na Tabela 3.1, que por sua vez foi superior ao desempenho na parte que avalia a Formação Geral, que teve média igual a 39,15, conforme mostrado na Tabela 3.2.
Os concluintes apresentaram um desempenho médio ( 54,66 ) superior ao dos ingressantes ( 47,80 ), havendo ainda indícios de maior homogeneidade entre as notas médias dos ingressantes do que entre as dos concluintes, como sugeriu o desvio padrão, menor no grupo dos ingressantes do que no dos concluintes.
As notas máximas obtidas por concluintes ( 100,0 ) e ingressantes ( 97,3 ) foram diferentes.
A distribuição de notas para o grupo dos concluintes se apresentaram com uma inclinação para notas maiores, como evidenciados pelo coeficiente de assimetria negativo.
A situação de assimetria é confirmada pelo Gráfico 3.5 que proporciona uma comparação do desempenho entre concluintes e ingressantes em relação ao Componente Específico da prova.
Claramente a distribuição das notas médias dos ingressantes, à direita daquela dos concluintes, evidenciou o percentual mais elevado de notas médias nas faixas mais baixas para este grupo, com exceção do primeiro, que como já há comentado, incluiu uma maior proporção de zeros entre os concluintes.
As medianas das distribuições das notas dos ingressantes e dos concluintes foram 48,2 e 55,5, respectivamente, enfatizando a superioridade dos concluintes neste parte da prova.
Assim como os Gráficos 3.2 e 3.4, o Gráfico 3.6 apresenta uma comparação dos resultados em relação às Regiões, categoria administrativa e Organização Acadêmica, agora levando em conta o desempenho dos ingressantes e dos concluintes no Componente Específico da prova.
Tendo como foco as notas médias dos estudantes, em determinada Região, observa -se que, entre os concluintes, as mais elevadas foram encontradas nas Regiões Nordeste ( 55,19 ), Sudeste ( 55,34 ) e Sul ( 55,38 ), valores acima da média nacional para este componente ( 54,66 ).
Entre os ingressantes, as maiores notas ocorreram nas Regiões : Sudeste ( 48,16 ) e Sul ( 50,40 ), todos acima da média nacional ( 47,79 ).
Novamente, para a análise segundo a categoria administrativa, somente entidades municipais e particulares foram representadas.
Observa -se comportamento diverso àquele da Formação Geral, as médias mais elevadas entre os concluintes ( 55,71 ) \textbf{estão} nas instituições municipais quando na parte relativa à Formação Geral \textbf{estavam} nas particulares e as médias mais elevadas entre os ingressantes ( 47,91 ) ficam também nas instituições municipais tanto para a parte relativa à Formação Geral quanto ao Componente Específico.
Quanto à Organização Acadêmica, o comentário é o mesmo já feito para a nota bruta e referente à Formação Geral : Universidades, Centros Universitários e Faculdades são representados neste curso.
Considerando -se o tipo de Organização Acadêmica, entre os concluintes, as Faculdades apresentaram os maiores valores, seguidas dos Centros Universitários e das Universidades.
E entre os ingressantes, Centros Universitários apresentaram os maiores valores, seguidos das Faculdades e das Universidades. 3.2 Análise das Questões Objetivas 3.2.1 Formação Geral A Tabela 3.4 apresenta as estatísticas básicas relativas às oito questões objetivas da Formação Geral dos estudantes.
Adiantamos que a primeira questão foi desconsiderada pelo critério da correlação ponto-bisserial.
Como pode ser observado na Tabela 3.4, os estudantes concluintes obtiveram desempenho superior, com nota média de 50,67, em relação aos estudantes ingressantes cuja nota média foi 47,63.
A variabilidade entre os concluintes foi um pouco menor do que entre os ingressantes, como indicaram os desvios-padrão de 21,6, para o grupo de concluintes e 21,7 para o grupo dos ingressantes.
A amplitude total das notas foi a mesma para ambos os grupos : nota mínima de 0,00 e máxima de 100,0.
A mediana foi superior para os concluintes.
A assimetria, medida pelo coeficiente de assimetria ( skewness ) foi nula indicando simetria para os grupos.
A Tabela 3.5 apresenta a classificação das questões objetivas da Formação Geral, segundo o índice de facilidade ( questões em vermelho foram desconsideradas pelo critério ponto-bisserial ).
Como as questões foram comuns a todas as áreas, a tabela apresenta a classificação considerando -se os concluintes de todas as áreas e aqueles do curso de TECNOLOGIA EM MARKETING.
De acordo com os índices obtidos, as questões objetivas da prova foram assim avaliadas para a área : das oito questões, duas tiveram entre 60 \%, inclusive, e 85 \% de acertos e foram classificadas como fáceis ( somente uma para o total de concluintes ) ; outras duas foram consideradas medianas, alcançando entre 40 \%, inclusive, e 60 \% de acertos ( três para o total de concluintes ) ; três obtiveram entre 15 \%, inclusive, e 40 \% de acertos, sendo consideradas difíceis ( três para o total de concluintes ).
Por fim, uma foi classificada como muito difícil tento obtido percentual menor do que 15 \% de acertos ( uma para o total de concluintes ), coincidentemente a questão desconsiderada pelo critério ponto-bisserial.
Nenhuma questão foi classificada na categoria muito fácil.
A Tabela 3.6 descreve os resultados das análises das questões objetivas relativas à Formação Geral, segundo o poder de discriminação, utilizando -se para tal o índice de discriminação ponto-bisserial.
Como as questões foram comuns a todas as áreas, a tabela apresenta a classificação considerando -se os concluintes de todas as áreas e aqueles do curso de TECNOLOGIA EM MARKETING.
A desconsideração da questão segundo este critério considerará o universo de todos os concluintes.
Nenhuma das oito questões apresentou índices iguais ou acima de 0,40 ( nenhuma para a população de concluintes como um todo ), sendo, portanto, classificadas como muito boas nesta característica, para este grupo de estudantes, uma foi classificada como boa com índice entre 0,30, inclusive, e 0,40 ( três para os concluintes como um todo ) e seis foram classificadas como médias, com índice entre 0,20, inclusive, e 0,30 ( quatro para os concluintes como um todo ).
Uma das questões foi classificada como fraca tanto para a população de concluintes como um todo como para a de concluintes de TECNOLOGIA EM MARKETING e foi desconsiderada pelo critério ponto-bisserial. de acertos, sendo consideradas difíceis ( três para o total de concluintes ).
Por fim, uma foi classificada como muito difícil tento obtido percentual menor do que 15 \% de acertos ( uma para o total de concluintes ), coincidentemente a questão desconsiderada pelo critério ponto-bisserial.
Nenhuma questão foi classificada na categoria muito fácil.
A Tabela 3.6 descreve os resultados das análises das questões objetivas relativas à Formação Geral, segundo o poder de discriminação, utilizando -se para tal o índice de discriminação ponto-bisserial.
Como as questões foram comuns a todas as áreas, a tabela apresenta a classificação considerando -se os concluintes de todas as áreas e aqueles do curso de TECNOLOGIA EM MARKETING.
A desconsideração da questão segundo este critério considerará o universo de todos os concluintes.
Nenhuma das oito questões apresentou índices iguais ou acima de 0,40 ( nenhuma para a população de concluintes como um todo ), sendo, portanto, classificadas como muito boas nesta característica, para este grupo de estudantes, uma foi classificada como boa com índice entre 0,30, inclusive, e 0,40 ( três para os concluintes como um todo ) e seis foram classificadas como médias, com índice entre 0,20, inclusive, e 0,30 ( quatro para os concluintes como um todo ).
Uma das questões foi classificada como fraca tanto para a população de concluintes como um todo como para a de concluintes de TECNOLOGIA EM MARKETING e foi desconsiderada pelo critério.
Na sequência, no Quadro 3.1, \textbf{estão} especificadas, por questão, as habilidades pelas questões da Formação Geral da prova da TECNOLOGIA EM MARKETING, bem como os resultados da análise relativa aos índices de facilidade e de discriminação de \textbf{cada} questão.
O índice de facilidade para área de TECNOLOGIA EM MARKETING teve uma variação considerável, oscilando entre 0,10 até 0,79 e o de discriminação ficou entre -0,01 até 0,38.
A questão 7 aferiu as habilidades de “ Analisar e criticar informações ”, “ Extrair conclusões por indução e/ou dedução ”, “ Estabelecer relações, comparações e contrastes em diferentes situações ”, “ Detectar contradições ”, “ Fazer escolhas valorativas avaliando consequências ”, “ Questionar a realidade ”, e “ Argumentar coerentemente ” foi a que apresentou, para a população de concluintes de TECNOLOGIA EM MARKETING, maior poder discriminatório com índice de 0,38.
A questão mais fácil, com 79 \% de acertos, foi a de número 8, que aferiu também, além das já listadas, a habilidade de “ Ler e interpretar textos ”.
De outra parte, a questão com menor índice de facilidade e de discriminação deste conjunto foi a questão 1, que foi desconsiderada, e que deveria aferir as habilidades de “ Ler e interpretar textos ”, “ Analisar e criticar informações ”, “ Extrair conclusões por indução e/ou dedução ”, “ Estabelecer relações, comparações e contrastes em diferentes situações ”, “ Detectar contradições ”, “ Fazer escolhas valorativas avaliando consequências ”, “ Questionar a realidade ”, e “ Argumentar coerentemente ”.
O Gráfico 3.7 mostra, a título de ilustração, o comportamento das respostas à questão 7 de Formação Geral os estudantes concluintes da área de TECNOLOGIA EM MARKETING como função da nota geral da prova4.
Trata -se de uma questão considerada fácil para os concluintes de TECNOLOGIA EM MARKETING e, de dificuldade média para o conjunto de concluintes de todas as áreas.
O poder de discriminação foi considerado bom tanto para os estudantes da área como para o conjunto de todos os concluintes.
Neste Gráfico, \textbf{cada} uma das cinco curvas, representou o percentual de respostas em determinada alternativa da questão, em função do escore total dos estudantes concluintes da área de TECNOLOGIA EM MARKETING.
A curva azul marinho corresponde à alternativa A, a correta para esta questão.
Assim, observa -se que entre os estudantes com notas mais baixas – desempenho geral mais fraco na prova como um todo –, houve a escolha de alternativas incorretas, especialmente a alternativa B – curva verde claro, além da alternativa correta.
A proporção de estudantes que selecionaram a resposta correta A, aumentou com o desempenho deste grupo de estudantes chegando ao acerto de 100 \% dos estudantes com notas pelo menos 87,5.
Esta análise permitiu verificar como a questão discriminou entre os dois grupos de desempenho.
Os Gráficos relativos às demais questões de Formação Geral constam do Anexo I. 3.2.2 Componente Específico A Tabela 3.7 apresenta as estatísticas básicas em relação às questões objetivas do Componente Específico.
Como podem ser observados, os estudantes concluintes obtiveram desempenho superior, com uma média de 58,20 pontos, em relação aos estudantes ingressantes com uma média de 51,36 pontos.
Os desvios-padrão indicaram que a variabilidade entre os concluintes foi maior ( 17,5 ) do que entre os ingressantes ( 17,3 ).
Os coeficientes de assimetria foram negativos indicando uma inclinação para à direita.
A análise das questões objetivas do Componente Específico em TECNOLOGIA EM MARKETING, em relação ao grau de facilidade da prova ( ver Tabela 3.8 ), mostrou que, das 27 questões, oito encontraram -se na categoria difícil ( entre 15 \% e 40 \%, não incluindo, de acertos ), onze são de nível médio ( entre 40 \% e 60 \%, não incluso, de acertos ), seis foram na categoria fácil ( entre 60 \% e 85 \%, não incluso, de acertos ) e nenhuma teve a classificação muito fácil ( acima de 85 \% de acertos ).
Duas questões tiveram percentual de acertos menor do que 15 \% ( classificação ' muito difícil ' ).
Nenhuma questão foi desconsiderada pela Comissão Assessora da Área de TECNOLOGIA EM MARKETING.
Das 27 questões objetivas relativas ao Componente Específico ( ver Tabela 3.9 ), seis obtiveram o índice de discriminação muito bom, sete obtiveram o índice de discriminação bom e seis obtiveram nível de discriminação médio.
O índice fraco que eliminaria a questão foi alcançado por oito questões.
Na sequência, no Quadro 3.2 \textbf{estão} especificadas, por questão, as habilidades aferidas pelas questões do Componente Específico da prova, bem como os resultados da análise relativa aos índices de facilidade e de discriminação de \textbf{cada} questão.
O Gráfico 3.8 mostra, a título de ilustração, o comportamento da questão 11 do Componente Específico para os estudantes concluintes da área de TECNOLOGIA EM MARKETING como função da nota geral na prova5.
Trata -se de uma questão considerada de dificuldade média para os concluintes de TECNOLOGIA EM MARKETING.
O poder de discriminação foi considerado bom para os estudantes da área.
Neste Gráfico, \textbf{cada} uma das cinco curvas, representa o percentual de respostas em determinada alternativa da questão, em função da nota total dos estudantes concluintes de TECNOLOGIA EM MARKETING.
A curva azul marinho correspondeu à alternativa A, a correta para esta questão.
Assim, observou -se que entre os estudantes com escores mais baixos – desempenho geral mais fraco na prova como um todo –, houve a escolha de alternativas incorretas, entre elas as alternativas E – curva em roxo e a alternativa D – curva em verde escuro.
A proporção de estudantes que selecionaram a resposta correta aumentou com o desempenho deste grupo de estudantes chegando ao acerto de aproximadamente 95 \% dos estudantes com notas pelo menos 87,5.
Esta análise permitiu verificar como a questão discriminou entre os grupos de desempenho.
Os Gráficos relativos às demais questões do Componente Específico constam do Anexo I. 3.3 Análise das Questões Discursivas 3.3.1 Formação Geral A Tabela 3.10 apresenta as estatísticas básicas da prova por grupo de estudantes.
A análise dos resultados de desempenho dos estudantes nas questões discursivas que são apresentados na tabela 3.10 e no gráfico 3.9, mostrou que as notas foram mais baixas no conjunto dessas questões que no das questões objetivas.
O total de estudantes inscritos para a prova foi de 23301.
Destes, 67,0 \% compareceram, sendo que a presença foi maior entre concluintes ( 78,1 \% ) do que entre ingressantes ( 59,5 \% ).
A média geral da prova foi de 24,42, sendo que os ingressantes obtiveram média mais baixa, de 22,54, que os concluintes, cuja média foi 26,56.
O desvio-padrão geral foi de 32,3, sendo o do grupo de ingressantes menor ( 31,5 ), indicando que, em relação ao grupo dos concluintes ( 33,2 ), aqueles possuíram uma distribuição mais homogênea nas notas.
A nota máxima foi 100,0, atingida por concluintes e ingressantes.
Os coeficientes de assimetria para os dois grupos considerados, o de concluintes e o de ingressantes, bem como para a população como um todo foram todos positivos, indicando uma maior concentração à esquerda da distribuição.
Para todos os grupos mais da metade dos estudantes presentes obtiveram nota zero nas questões discursivas da Formação Geral.
Das 9.164 provas com nota zero, 330 ( 3,6 \% ) continham um protesto, 1.789 ( 19,5 \% ) \textbf{estavam} em branco e 7.045 ( 76,9 \% ) foram respondidas, porém de forma inadequada.
Entre os concluintes a maioria ( 79,1 \% ) dentre aqueles com nota zero apresentaram as provas respondidas de forma inadequada.
Entre os ingressantes 75,1 \% das notas nulas ( zero ) decorreram do desempenho do próprio estudante verificadas na correção.
Uma comparação entre o desempenho de ingressantes e concluintes pode ser realizada a partir dos dados apresentados no Gráfico 3.9.
As notas dos concluintes, de um modo geral, foram ligeiramente mais elevadas do que as dos ingressantes, cuja distribuição de notas \textbf{está} à direita dos primeiros, superioridade reforçada pela proporção de notas nulas : 55,6 \% entre os concluintes e 61,5 \% entre os ingressantes.
Na sequência, discorrer -se-á, especificamente, sobre uma das questões discursivas de Formação Geral, a número 10 já que a número 9 foi desconsiderada.
Em especial, serão analisados os conteúdos englobados na questão e os resultados obtidos pelos estudantes.
Ao final, são apresentados os comentários da banca corretora a respeito do observado nas respostas dos estudantes, suas impressões e conclusões.
Cumpre esclarecer que, tendo em vista que as questões discursivas de Formação Geral foram padronizadas, ou seja, constam de todas as provas, os comentários da banca foram os mesmos para todas as áreas acadêmicas, eis que direcionados a todos os estudantes que participaram do ENADE/2009. 3.3.1.1 Análise da Questão Discursiva 9 de Formação Geral DESCONSIDERADA 3.3.1.2 Análise da Questão Discursiva 10 de Formação Geral A Questão Discursiva nº 10 buscou avaliar as habilidades : • Ler e interpretar textos ; • Analisar e criticar informações ; • Extrair conclusões por indução e/ou dedução ; • Estabelecer relações, comparações e contrastes em diferentes situações ; • Fazer escolhas valorativas avaliando consequências ; • Questionar a realidade ; • Argumentar coerentemente.
O Quadro 3.3 apresenta as habilidades/competências, temas predominantes e estatísticas básicas da questão 10 de Formação Geral por grupo de estudantes.
Cumpre notar que como a questão 9 foi desconsiderada a nota referente às questões discursivas na Formação Geral é idêntica a nota da questão 10.
A média geral da prova foi de 24,42, sendo que os ingressantes obtiveram média mais baixa, de 22,54, que os concluintes, cuja média foi 26,66.
O desvio-padrão geral foi de 32,3, sendo o do grupo de ingressantes menor ( 31,5 ), indicando que, em relação ao grupo dos concluintes ( 33,2 ), aqueles possuíram uma distribuição mais homogênea nas notas.
A nota máxima foi 100,0, atingida por concluintes e ingressantes.
Os coeficientes de assimetria para os dois grupos considerados, o de concluintes e o de ingressantes, bem como para a população como um todo foram positivos, indicando uma maior concentração à esquerda da distribuição.
Para todos os grupos, mais da metade dos estudantes presentes tiveram nota zero nesta questão discursiva da Formação Geral.
Das 9.164 provas com nota zero, 330 ( 3,6 \% ) continham um protesto, 1.789 ( 19,5 \% ) \textbf{estavam} em branco e 7.045 ( 76,9 \% ) foram respondidas, porém de forma inadequada.
Entre os concluintes a maioria ( 79,1 \% ) dentre aqueles com nota zero \textbf{estavam} com as provas respondidas de forma inadequada.
Entre os ingressantes, 75,1 \% das notas nulas ( zero ) decorreram do desempenho do próprio estudante verificadas na correção.
Uma comparação entre o desempenho de ingressantes e concluintes para esta questão foi exatamente igual à já realizada para a nota discursiva de Formação Geral a partir dos dados apresentados no Gráfico 3.9. 3.3.1.3 Comentários sobre as respostas à Questão Discursiva 10 Apresentação A questão 10 de Formação Geral apresentou um questionamento capaz de analisar a habilidade do estudante em reconhecer o que é um direito social, optar por um direito social vigente que possa diminuir as desigualdades de renda familiar no Brasil e discutir os argumentos em favor do direito social escolhido como suporte à resposta.
Para análise das respostas apresentadas pelos estudantes, houve uma leitura inicial bastante criteriosa.
Foram considerados para efeito de pontuação, os seguintes argumentos : • Acesso à educação pública, gratuita e de qualidade, garantindo ao cidadão melhores condições de vida ; • Assistência à saúde, em seu contexto mais amplo, garantindo os cuidados essenciais de saúde, baseados em técnicas apropriadas ao alcance de todas as pessoas, contando com a participação da população ; • Permanência do estudante na escola, em todos os níveis escolares ( da educação infantil a educação superior, possibilitando a qualificação profissional ), o que garante melhores condições de trabalho e remuneração.
Considerou -se a qualificação e o desempenho docente, desenvolvendo nos estudantes melhores expectativas de aprendizagem e interesse pelo conteúdo da escola.
Os professores desenvolvendo metodologias adequadas e conhecimento pleno quanto a questões de aprendizagem ; • Condições dignas de trabalho, com remuneração que garanta qualidade de vida.
Foram consideradas as condições de trabalho, a preparação profissional para o trabalho e os cursos de aperfeiçoamento, oportunizando momentos de aprendizagem e desenvolvimento, a partir de situações desencadeadoras ; • Ser proprietário do imóvel que reside, reduzindo os gastos com aluguel, garantindo melhores condições de vida, equilíbrio financeiro e construção do bem-estar familiar.
Desempenho dos estudantes na questão É importante salientar que havia respostas que desenvolveram adequadamente o questionamento apresentado, com fundamentação baseada em argumentos coerentes.
Contudo, em um número considerável de respostas, havia apenas um argumento e não dois, conforme fora solicitado.
Muitos estudantes informaram sobre problemas da educação, de um modo geral.
Algumas observações mencionadas pelos corretores durante o processo de correção foram : • há preocupação com uma educação de qualidade, com sugestões de cursos profissionalizantes e métodos preparatórios de aulas. • a educação é revelada como de baixa qualidade, exaltando -se treinamentos e cursos de atualização para a classe docente. • a saúde como essência para a vida.
O SUS é colocado como parte integrante de direito social, garantindo às pessoas fácil acesso. • oportunidade de emprego como forma de ascensão social, direito, cidadania. • os políticos brasileiros como causadores do despreparo educacional.
As respostas analisadas apresentaram norma culta, coesão e coerência razoáveis, que garantiram a compreensão da mensagem desenvolvida pelos estudantes.
Os erros ortográficos foram poucos e não comprometeram a discussão sobre o assunto e sobre os pontos de vista abordados.
Quanto à coerência, houve um desvio em relação ao objetivo da questão, já que alguns estudantes, através dos seus argumentos retrataram reivindicações, reclamações e situações de problemas enfrentados pelos brasileiros, citando golpes políticos.
Houve erros quanto à concordância verbal e nominal.
Houve, também, dificuldade em utilizar os conectores corretamente.
Alguns estudantes apresentaram em suas respostas redundância e prolixidade.
O desempenho foi considerado satisfatório, pois a maioria dos estudantes soube responder ao questionamento.
Houve domínio do conteúdo, pois as ideias, os argumentos e a originalidade corresponderam ao esperado. 3.3.2 Componente Específico Nas questões discursivas do Componente Específico – questões de números 38, 39 e 40 –, a diferença de desempenho entre os grupos não foi muito acentuada, inclusive menor do que para as questões discursivas de Formação Geral.
A Tabela 3.11 contém as estatísticas da média destas questões.
Observa -se que no grupo dos concluintes a nota média alcançada foi 34,93, enquanto que no grupo dos ingressantes foi de 27,92, ambas maiores do que as obtidas nas questões de Formação Geral, que foram de 22,5 para os ingressantes e 26,6 para os concluintes ( Tabela 3.10 ).
As notas são, porém, menores dos que as obtidas para as questões objetivas desse componente ( Tabela 3.7 ) : respectivamente 51,4 e 58,2 para ingressantes e concluintes.
A análise de \textbf{cada} uma destas questões será feita nas subseções seguintes.
Os valores positivos para a medida de assimetria indicaram uma distribuição com cauda para a direita.
Dos 23301 presentes, 25,2 \% obtiveram nota zero.
A tabela 3.12 apresenta um \textbf{detalhamento} da situação destes estudantes.
A maioria ( 45,8 \% ) não fez nenhuma das 3 questões.
Esta proporção foi ligeiramente maior entre os ingressantes ( 48,3 \% ) do que entre os concluintes ( 42,0 \% ).
A combinação de questões em branco e preenchidas correspondeu a um pouco menos de 10 \% das provas. 0,3 \% no total das provas correspondeu a 3 questões com protesto.
O Gráfico 3.10 apresenta a distribuição das notas nas questões discursivas do Componente Específico por grupo de estudantes.
A moda se situou no primeiro intervalo, aquele com valores entre 0 e 10 incluindo os extremos do intervalo.
A impressão de assimetria informada pelos coeficientes de assimetria positivos ( Tabela 3.11 ) foi reforçada no gráfico pela concentração no primeiro intervalo. 3.3.2.1 Análise da Questão Discursiva 38 do Componente Específico A questão 38, cujas habilidades aferidas encontram -se descritas no Quadro 3.4, que ainda contém outras informações, seguido da Tabela 3.13 e do Gráfico 3.11, foi aquela em que os grupos de estudantes estudados observados obtiveram o melhor desempenho entre as questões discursivas do Componente Específico.
A nota média do grupo dos ingressantes foi de 37,2, e a dos concluintes 46,7.
O percentual de notas zero obtidas nos dois grupos foi o menor das questões do Componente Específico com 38,9 \%, maior entre os ingressantes ( 43,4 \% ) do que entre os concluintes ( 33,7 \% ).
Quase 2/5 dos concluintes e dos ingressantes que tiraram zero deixaram a questão em branco ( ver Tabela 3.13 ).
Uma pequena minoria ( 4,2 \% ) utilizou a questão para protestar.
Apesar do baixo desempenho, em ambos os grupos a nota máxima de 100,0 pontos foi alcançada.
Os coeficientes de assimetria positivos indicam uma inclinação para a esquerda e cauda para a direita.
O Gráfico 3.11 apresenta a distribuição das notas na questão 38 discursiva do Componente Específico por grupo de estudantes.
A moda se situou no primeiro intervalo, aquele com valores entre 0 e 10 incluindo os extremos do intervalo.
Nota -se, ainda, dois máximos locais, um no sexto intervalo ( notas entre 50 e 60 ) e outro no último intervalo de notas maiores.
A impressão de assimetria informada pelos coeficientes de assimetria positivos ( Quadro 3.5 ) foi reforçada no gráfico pela concentração no primeiro intervalo, porém, os outros máximos locais sugerem \textbf{certa} simetria além de uma polarização nas notas dos estudantes. 3.3.2.2 Comentários sobre as respostas à Questão Discursiva 38 O objetivo dessa questão era avaliar o conhecimento dos estudantes sobre extensão de marca ( foi solicitado que fosse uma questão sobre produto e extensão de marca, é um aspecto importante da gestão de produtos ).
A banca de elaboração de questões decidiu por elaborar questões com complexidade crescente.
Portanto, consideramos que esta questão era, no nível de dificuldade, considerada fácil/média.
Primeiramente, destacamos que a questão foi de baixa/média complexidade, não demandando conhecimento profundo sobre o assunto, apenas que se conhece o conceito de extensão de linha de produto, sendo este diferente do conceito de extensão de categoria.
Foi solicitado na questão que os estudantes sugerissem duas extensões de linha possíveis.
Muitos estudantes responderam corretamente.
Consignamos métodos dos pontos aos estudantes que citaram duas extensões de categoria.
Entretanto, vários estudantes não souberam responder.
Observou isto tanto pela não resposta ( ou resposta em branco ), quanto respondendo completamente errado.
Portanto, alguns estudantes aparentaram não compreender o solicitado na questão.
Outro erro comum foi a citação de uma extensão de linha e a indicação de uma ferramenta do composto promocional ( por exemplo : linha light e propaganda ).
Apesar destes erros, que acreditamos que tenham sido cometidos pelos ingressantes, a impressão geral foi de que os estudantes possuem conhecimento sobre extensão de linha, havendo uma confusão entre este conceito e o de extensão de categoria.
Coesão e coerência na resposta não foram questões importantes, visto que bastava ao estudante citar as duas extensões de marca, entretanto, foi recomendado aos corretores que só dessem os pontos inteiros de coesão e coerência para aqueles que respondessem com uma frase completa e retirassem pontos por português errado ( como no caso de LAITE - light ).
Concluindo, percebeu -se que a grande maioria dos estudantes que \textbf{estão} terminando o curso acertou esta questão. 3.3.2.3 Análise da Questão Discursiva 39 do Componente Específico A questão 39, cujas habilidades aferidas encontram -se descritas no Quadro 3.5, que ainda contém outras informações, seguido da Tabela 3.14 e do Gráfico 3.12, foi aquela em que os grupos de estudantes estudados observados obtiveram o desempenho médio igual a outra questão discursiva do Componente Específico ( 40 ), com média geral em 26,0.
A nota média do grupo dos ingressantes foi de 22,6, e a dos concluintes 29,8.
O percentual de notas zero obtidas nos dois grupos foi de 52,0 \% entre os ingressantes e 42,2 \% entre os concluintes.
Das provas com notas nulas nesta questão, cerca de 2/5 foi deixada em branco : 42,5 \% entre os ingressantes e 33,9 \% entre os concluintes ( ver Tabela 3.14 ).
Entre os com nota nula, uma pequena minoria ( 4,0 \% ) utilizou a questão para protestar.
Apesar do baixo desempenho, em ambos os grupos a nota máxima de 100,0 pontos foi alcançada.
Os coeficientes de assimetria positivos indicaram uma inclinação para a esquerda e cauda para a direita.
O Gráfico 3.12 apresenta a distribuição das notas na questão 39 discursiva do Componente Específico por grupo de estudantes.
A moda se situou no primeiro intervalo, aquele com valores entre 0 e 10 incluindo os extremos do intervalo.
A impressão de assimetria informada pelos coeficientes de assimetria positivos ( Quadro 3.5 ) foi reforçada no gráfico pela concentração no primeiro intervalo.
A distribuição apresentou um máximo local no sexto intervalo, de notas medianas. 3.3.2.4 Comentários sobre as respostas à Questão Discursiva 39 O objetivo dessa questão era verificar o conhecimento sobre composto promocional.
Conforme dito anteriormente, a banca de elaboração de questões decidiu por elaborar questões com complexidade crescente.
Portanto, consideramos que esta questão era, no nível de dificuldade, considerada média.
Assim, a questão não demandava conhecimento profundo sobre composto promocional ( item pedido para esta questão ), mas que conhecessem o conceito de composto promocional, suas ferramentas e o motivo de empregá -las.
Consideramos que a resposta \textbf{estaria} corretamente respondida se o estudante citasse duas ferramentas do composto promocional, além da sua justificativa de uso ( como melhor visibilidade da marca, aumento de vendas, melhorar a percepção do consumidor sobre a marca, etc. ).
Sobre o padrão de respostas encontrado o que foi observado é que a maioria dos estudantes conseguiu citar ferramentas do composto promocional pertinentes.
Entretanto, encontramos problemas não só na elaboração de justificativas coerentes, mas também na falta de coerência do próprio uso da língua portuguesa.
Nesta questão, notou -se a deficiência na expressão escrita dos estudantes.
Outro erro comum foi citar a mesma ferramenta duas vezes, indicando não compreenderem o solicitado ou não perceberem a semelhança entre as ferramentas ( como promoção para crianças e para homens ).
Alguns estudantes confundiram o solicitado, como nos casos em que ou citaram slogans de empresas existentes ( “ Nestlé, todo mundo conhece, todo mundo confia. ” ), ou mesmo criaram algum slogan ( “ Compre Smack e se livre do stress. ” ).
Outros erros comuns foram a confusão entre composto promocional e composto de marketing, quando alguns estudantes citaram os outros 3 P`s ( preço, praça e até produto, novamente ).
Houve falta de percepção em relação aos termos propaganda e publicidade, que são usados de forma diferente e não como sinônimos, assim alguns responderam “ publicidade em rádio e TV ”.
Embora muitos estudantes aplicassem corretamente o conceito de merchandising, ainda assim encontramos algumas respostas falando em “ merchandising em programa de TV ”.
É importante salientar que os estudantes demonstraram não compreenderem não só as características das ferramentas, mas principalmente, sua utilidade.
Salienta -se, também, que a deficiência encontrada na expressão escrita ( seja na falta de uso da norma culta, seja na construção de argumentos para responder à questão ) pode impactar outras habilidades gerenciais, visto que pode vir a impactar nas habilidades de análise, planejamento, implementação e controle como por exemplo : ( “... atrair os clientes impotenciais ” ) ( “... ficando recioso em comprar ” ). 3.3.2.5 Análise da Questão Discursiva 40 do Componente Específico A questão 40, cujas habilidades aferidas encontram -se descritas no Quadro 3.6, que ainda contém outras informações, seguido da Tabela 3.15 e do Gráfico 3.13.
A questão foi aquela em que os grupos de estudantes observados obtiveram o desempenho médio igual a outra questão discursiva do Componente Específico ( 39 ), com média geral em 26,0.
A nota média do grupo dos ingressantes foi de 24,0, e a dos concluintes 28,4.
O percentual de notas nulas : 49,5 \% entre os ingressantes e 43,0 \% entre os concluintes.
Aqui, a maioria das provas com nota zero decorreu de provas com respostas mal formuladas : 65,0 \% entre os ingressantes e 73,2 \% entre os concluintes ( ver Tabela 3.15 ).
Uma pequena minoria ( 0,6 \% ), dentre aqueles com nota nula, utilizou a questão para protestar.
Em ambos os grupos a nota máxima de 100,0 pontos foi alcançada.
Os coeficientes de assimetria positivos indicaram uma inclinação para a esquerda.
O Gráfico 3.13 apresenta a distribuição das notas na questão 40 discursiva do Componente Específico por grupo de estudantes.
A moda se situou no primeiro intervalo, aquele com valores entre 0 e 10 incluindo os extremos do intervalo.
A impressão de assimetria informada pelos coeficientes de assimetria positivos ( Quadro 3.6 ) foi reforçada no gráfico pela concentração no primeiro intervalo. 3.3.2.6 Comentários sobre as respostas à Questão Discursiva 40 O objetivo dessa questão era verificar o conhecimento sobre comportamento do consumidor, conforme solicitado pela Comissão Assessora nomeada pela DAES/INEP/MEC, que também sugeriu que fosse de grau difícil de resposta.
Esta questão exigia do estudante o conhecimento sobre comportamento do consumidor, exigindo que ele não só conhecesse o conceito de fatores pessoais, mas também articulasse sua resposta com base no texto.
Encontramos nesta questão o maior número de respostas erradas.
Na maioria dos casos, os estudantes não compreenderam o que foi solicitado.
O estudante deveria ser capaz de identificar no texto elementos que evidenciassem como os fatores pessoais poderiam influenciar as decisões dos consumidores, concordando ou discordando da decisão do gerente.
O aparente motivo para o erro foi na articulação da elaboração da resposta.
Muitos estudantes se valeram de “ achismos ” sem nenhuma base teórica ou do texto como por exemplo : ( “ chocolate é o mais gostoso dos afrodisíacos ; ingerido com moderação, chocolate é o máximo para a sexualidade do homem ” ).
Acreditamos que a falta de preocupação com fatos e o uso extensivo de preconceitos na formulação das respostas pode levar, na vida profissional, problemas na formulação de estratégias de marketing.
Esta preocupação também foi reforçada quando os estudantes afirmaram que “ todo mundo come chocolate ” que não compreendem nem sabem aplicar os conceitos de segmentação e posicionamento.
Nesta questão, se fez notar de maneira especial a deficiência na expressão escrita, de modo geral.
Encontramos, em muitas respostas dos estudantes, a falta de articulação da linguagem para a construção de argumentos, visando à resposta da questão.
De modo geral, os estudantes apresentaram flagrante incompreensão do conceito de fatores pessoais e insuficiente entendimento do que foi pedido na questão, o que torna indispensável reforçar o ensino de comportamento do consumidor e sua articulação com estratégia de marketing.

% matched lemmas: cada, certo, detalhamento, estar
\end{textsample}
